\chapter{物种怎样分开}

\begin{kaichang}
动物园的海报上写着：“狮虎兽——狮子和老虎的孩子！”旁边还贴着一张照片：一头毛色浅黄、带着淡淡条纹的大家伙，比它的父母都大一圈。

等等。狮子和老虎不是两个不同的物种吗？中学生物课上说，不同物种之间不能生出能育的后代——那这只狮虎兽是怎么来的？它算是狮子，还是老虎，还是一个新物种？反过来问：你家的狗和郊外的狼，长得天差地别，却能生出健康的小狗，它们算一个物种还是两个？

先别急着查答案。把你现在的猜测写下来。这一章结束时你会发现，“物种”这两个字，比课本上的定义要麻烦得多——而这份麻烦，恰恰是理解进化历史最好的入口。
\end{kaichang}

\section{“物种”两个字，没你想的那么清楚}

先从能看见的现象说起。

马和驴交配，能生下一头骡子。骡子力气大、耐粗饲，农民很喜欢它——但骡子几乎从不当爸爸妈妈，它不育。狮子和老虎在人工饲养下偶尔生下狮虎兽，雄狮虎兽同样不育。狗呢？吉娃娃和藏獒外形差得像两种动物，可它们和灰狼之间都能生出能育的后代。再看植物：卷心菜、花椰菜、西兰花、芥蓝，摆在一起你绝不会认错——但它们全是同一个物种（甘蓝）的不同栽培品种，互相之间完全可以授粉结实。

人类还会故意利用这条界线。水果摊上的香蕉几乎找不到种子，无子西瓜也没有正经的瓜籽——它们都是三倍体植物，染色体配对在减数分裂时乱成一团，种子发育中途夭折。农民求之不得的“无子”，本质上就是一次人为制造的合子后隔离。

于是你看到一组古怪的反差：有的“长得像一家”的不能生，有的“长得不像一家”的偏能生。可见“物种”的界线既不画在外表上，也不是一道干脆的闸门。要弄明白这条线画在哪，得先把镜头推进细胞里，看看“生孩子”这件事在分子层面到底需要什么。

\section{生不出孩子，问题出在染色体上}

回忆第 69 章讲过的减数分裂：制造精子或卵细胞时，细胞里的染色体要两两配对——来自父亲的那一套和来自母亲的那一套，像两队人手拉手站好，再彼此交换片段、平均分配。配对能进行，前提是两条染色体足够相似：长度相近、基因排列的顺序相近，像两本页码一致的说明书才能逐页对照。

现在看骡子的处境。马的体细胞有 64 条染色体（32 对），驴有 62 条（31 对）。骡子从马那里得到 32 条，从驴那里得到 31 条，总共 63 条。到了骡子自己制造生殖细胞时，问题来了：马的那 32 条和驴的那 31 条，数目对不上，序列也差得太多，很多染色体找不到能配对的伙伴。配对一乱，减数分裂就崩盘，造不出正常的精子或卵子。骡子这头“混合动力车”本身跑得挺好，却没法把这份动力传下去。

狮虎兽的不育，道理更深一层。狮子和老虎的染色体数目倒是相同，但几百万年各自生活，两套基因组的“用词习惯”已经悄悄变了：同样的基因，在狮子里什么时候开、开多大音量，和在老虎里并不一样。两套基因挤在一个细胞里，有些开关指令互相打架，生殖细胞的发育就出故障（雄性受害尤其重——这背后有个规律，本章挑战层会提一句）。

把镜头再拉远一步，你已经能看到物种分开的粒子层图景了：\textbf{所谓两个物种，就是两套曾经相同的 DNA 说明书，被改到再也对不上页码的程度}。接下来的问题是：它们是怎么被改开的？

\section{到底什么算一个物种}

“改到什么程度才算两个物种？”这个问题，生物学家自己至今也没有唯一答案。他们手上有不止一把尺子，各有各的用法和短板：

\begin{table}[htbp]
\centering
\begin{tabular}{p{3.2cm}p{4.6cm}p{4.6cm}}
\toprule
物种概念 & 判断标准 & 明显的短板 \\
\midrule
生物学物种概念 & 自然状态下能否交配并产生可育后代 & 管不了不交配繁殖的生物（细菌、许多植物），也管不了化石 \\
形态学物种概念 & 外形和结构是否足够相似 & “像不像”有时靠主观，同种生物雌雄、幼成体可以差得很远 \\
系统发育物种概念 & 是否在进化树上构成独立的一支 & 需要先重建进化树，而重建本身依赖其他证据 \\
\bottomrule
\end{tabular}
\caption{三种常用的“物种”尺子。没有哪一把在所有场合都好用，生物学家按问题选尺子。}
\end{table}

中学课本讲的“生殖隔离”，属于第一把尺子——生物学物种概念。它对大多数动物很好用，也是本章的主力工具。但请你记住表格最后一列：这把尺子量不了细菌（它们不靠两性生殖），量不了化石（没法让两具骨架交配试试），遇到环物种还会打结（本章后面会讲）。所以严谨的说法是：\textbf{“物种”是人类给连续的自然界画的一套格子，格子非常有用，但自然界没有义务恰好落进格子里。}

\section{两个群体怎样越走越远}

既然尺子选好了，就用它来看物种分开的实际过程。

设想一座山，山谷里住着一群同一种的小鸟。某年地质变动，一条大河切开山谷，把鸟群分成东西两半，河水宽到它们飞不过去。这一步叫\textbf{地理隔离}——它只是按下了计时器，本身还不算物种形成。接下来，两个群体各自过各自的日子：第 76 章的自然选择、第 77 章的突变、遗传漂变，在两边独立地进行。西岸干燥，喙粗一点的个体吃硬种子占便宜；东岸湿润，喙细一点的个体啄虫子更顺手。几千年、几万代之后，两群鸟的喙形、羽色、求偶的歌声，都越差越远。

终于有一天，河道改道，两群鸟又碰面了。这时会出现三种结局：

\begin{itemize}[itemsep=2pt]
\item 它们照常互相求偶、交配，后代健康能育——隔离时间太短，差异不够，还算一个物种；
\item 它们互相看不上：歌声对不上暗号，或者繁殖季节错开了，根本不到交配那一步——这叫\textbf{合子前隔离}；
\item 它们能交配，但胚胎发育到一半停止，或者后代像骡子一样不育——这叫\textbf{合子后隔离}。
\end{itemize}

后两种结局合称\textbf{生殖隔离}。一旦出现稳定的生殖隔离，两群鸟就算真正分成了两个物种——它们各自的基因变化从此不再共享，像两条分岔后各自延伸的路。

地理隔离是最常见、也最好理解的开端，但不是唯一的。有些物种在同一片水域或同一座山里就分了家：比如一群果蝇里，偏爱在苹果上产卵的个体和偏爱在山楂上产卵的个体渐渐只在各自的果树上找对象，基因流自己断了线；植物更是能靠一次染色体加倍“瞬间成种”——三倍体不育，可若染色体再翻一倍成六倍体，配对又恢复正常，一株新植物当场就和祖先种生殖隔离了。物种形成的路径有很多条，地理隔离只是其中最常被写进课本的一条。

\begin{table}[htbp]
\centering
\begin{tabular}{p{2.6cm}p{3.4cm}p{6.2cm}}
\toprule
类型 & 机制 & 例子 \\
\midrule
合子前隔离 & 繁殖时间错开 & 两种松树在同一座山上，一个二月散粉，一个四月散粉 \\
 & 求偶行为不合 & 萤火虫靠闪光密码找对象，密码不对不予理睬 \\
 & 栖息微环境不同 & 同一湖里，一种鱼在浅水产卵，一种在深水产卵 \\
合子后隔离 & 杂种不能存活 & 两种蛙的受精卵发育到蝌蚪早期就停止 \\
 & 杂种不育 & 骡子、雄性狮虎兽 \\
\bottomrule
\end{tabular}
\caption{生殖隔离的常见机制。合子前隔离像“约不上会”，合子后隔离像“合作到一半散伙”。}
\end{table}

\begin{qiaoliang}
物种分开没有一声发令枪，差异是一点一点攒出来的——攒得有多快？算一笔账。第 77 章讲过，每个人出生时平均携带约 70 个父母都没有的新突变（二倍体基因组约 $6\times10^{9}$ 个碱基，乘以每代每位点约 $1.2\times10^{-8}$ 的突变率）。假设两个群体彻底隔离、每 20 年繁殖一代，那么一百万年里每个群体大约繁殖 5 万代，仅新突变一项就能积累约 $70\times 50000=350$ 万个差异位点。这还没算上自然选择会放大其中一部分差异的频率。也就是说，\textbf{即使什么都不“努力”，时间本身就在不停地改写着两套说明书}——地理隔离一旦维持百万年的量级，生殖隔离几乎是大概率事件。反过来，如果两群之间一直有少数个体跑来跑去通婚（基因流），差异就会被不断“冲淡”，物种就迟迟分不开。
\end{qiaoliang}

\section{把历史画成一棵树}

物种一旦分开，就不再合并（动物里极少例外）。于是，一部物种形成史天然是一棵不断分叉的树：一个祖先物种分成两支，每支再分成更细的枝。把这套关系画下来，就是\textbf{系统发育树}——它是符号层的主角，也是本章最容易被读错的图。

下面这棵树画的是人和几种现生灵长类的亲缘关系。注意，这不是它们的长相图，而是“分家顺序”的示意图：

\begin{center}
\begin{tikzpicture}[font=\small]
\draw (5.0,4.0) node[right]{人} -- (4.2,4.0);
\draw (5.0,3.2) node[right]{黑猩猩} -- (4.2,3.2);
\draw (4.2,3.2) -- (4.2,4.0);
\draw (5.0,2.4) node[right]{大猩猩} -- (3.4,2.4);
\draw (3.4,2.4) -- (3.4,3.6);
\draw (3.4,3.6) -- (4.2,3.6);
\draw (5.0,1.6) node[right]{红毛猩猩} -- (2.6,1.6);
\draw (2.6,1.6) -- (2.6,3.0);
\draw (2.6,3.0) -- (3.4,3.0);
\draw (5.0,0.8) node[right]{猕猴} -- (1.8,0.8);
\draw (1.8,0.8) -- (1.8,2.3);
\draw (1.8,2.3) -- (2.6,2.3);
\draw (1.0,1.55) -- (1.8,1.55);
\fill (4.2,3.6) circle (1.2pt) node[above left=1pt]{\footnotesize A};
\fill (1.8,1.55) circle (1.2pt) node[below left=1pt]{\footnotesize B};
\end{tikzpicture}
\end{center}

读这棵树只有三条规则：

\begin{itemize}[itemsep=2pt]
\item \textbf{分叉点是共同祖先。}节点 A 代表人和黑猩猩的最近共同祖先——一个生活在大约六七百万年前的古猿种群。它不是人，也不是今天的黑猩猩，而是两者共同的“外婆的外婆的……外婆”。
\item \textbf{树枝可以绕节点自由旋转。}把黑猩猩画在人上面，或把人画在最左边，树的意思一模一样。树梢的左右排列是画图人的排版习惯，不代表任何“先进程度”。
\item \textbf{所有树梢都是“现代人”。}树的末端全是今天活着的物种，它们在地球上存在的时间一样长——都从 38 亿年前的共同祖先一路传到了现在。猕猴不是“停在半路的猴子”，它只是和我们分家更早的亲戚，在自己那条路上同样进化了三千多万年。
\end{itemize}

\begin{wuqu}
“进化是从低等到高等的阶梯：细菌最原始，人是最高级；照这么进化下去，现在的猴子迟早会变成人。”——这是对系统发育树最深的误读。第一，树不是阶梯：阶梯意味着有方向、有顶端，而进化树的每次分叉只是两个群体各自适应各自的环境，不存在“谁在朝谁爬”。细菌在地球上活了 38 亿年，至今仍是数量最庞大、分布最广的生命，按“活得成功”这把尺子它们一点不输。第二，人不是从今天的猴子变来的：人和黑猩猩的关系是“堂兄弟”，不是“祖孙”——我们共享节点 A 那个祖先，之后各走各路。问“猴子为什么还没变成人”，就像问“你表哥为什么还没变成你”——他有自己的家谱和人生，从来就不在变成你的路上。把阶梯图从你的脑子里删掉，换成一棵不断分叉、没有顶端的灌木，是这一章最重要的一个动作。
\end{wuqu}

这套画法可以一直推到最深的地方。把现存所有生物的序列和结构证据叠在一起，它们能连成唯一一棵大树：细菌、古菌和包括你在内的真核生物是三大主枝，树根处站着所有生命的最后共同祖先——一个生活在大约 38 亿年前的古老细胞群体。你和窗台上那盆绿萝、培养皿里的大肠杆菌，都能在这棵树上找到彼此相连的分叉点。这不是比喻：你们细胞里合成蛋白质的核糖体、读取 DNA 的基本密码，用的是同一套零件和同一套语言（第 71 章），这正是“同源”最有力的注脚。至于病毒——它们没有自己的核糖体，算不算这棵树上的正式成员，学界还在争论，第 79 章会再碰这个话题。

那么，凭什么相信这棵树画对了？凭什么说鲸的祖先曾经在陆地上走路？下一节讲科学家怎样知道。

\begin{zhengju}
19 世纪的人看鲸，只觉得它们是“大鱼”。达尔文在《物种起源》里斗胆猜想过熊下水捕食可能演化成鲸一类的东西，被当时的报纸嘲笑了很久。谁也没想到，一百多年后，三种完全不同的证据会在这件事上汇合。

第一路证据来自\textbf{化石}。1979 年起，古生物学家在巴基斯坦的古老河床里挖出一连串生活于约五千万年前的动物：巴基鲸有典型的鲸式耳骨，却长着能走路的四肢；稍晚的走鲸四肢粗壮，既能在岸上撑住身体，又能划水；再晚的龙王鲸已经是全水生的庞然大物，却还挂着一对小得可笑的后腿——腿骨完整，有关节，只是再也够不着地面。一连串化石按年代排开，肢骨一步一步缩短，鼻孔一步一步往头顶挪。

第二路证据来自\textbf{比较解剖与胚胎}。今天鲸的鳍状肢里，骨骼排列和你手臂里的惊人相似：一根上臂骨、两根前臂骨、腕骨、五根指骨，只是指头被包在了鳍里。更妙的是，鲸的胚胎发育早期会长出后肢的芽，随后停止生长——发育过程把早已用不上的“祖传图纸”短暂地翻出来又收回去。

第三路证据来自\textbf{分子序列}。20 世纪 90 年代，科学家比较鲸和其他哺乳动物的 DNA 序列，结果让所有人愣了一下：和鲸序列最像的现生动物不是任何“鱼一样的家伙”，而是河马。随后人们在河马和鲸的基因组里找到多处相同的“旧伤疤”（插入在相同位置的病毒遗留序列），又在化石记录里找到了两者共同祖先的候选——一类叫石炭兽的古代偶蹄动物。鲸甚至被重新归进了偶蹄目：从分类上说，鲸是一种“下海的河马亲戚”。

三路证据来自完全不同的方法、不同年代、不同实验室，却指向同一棵树。科学上管这种情形叫“证据的汇聚”——一条结论如果被互不相干的证据链同时支持，它出错的概率就远远小于任何一条单独链出错的概率。当然，树形图的细节仍在被修正（比如某些早期鲸类的确切位置），这正是科学的工作方式：大方向靠汇聚的证据站稳，小细节留给新化石去改。
\end{zhengju}

\section{这棵树也有不好用的时候}

好的模型要知道自己的边界。“两个群体分开就再也合不上”这个图景，在不少场合会被自然界当面顶回来：

\begin{itemize}[itemsep=2pt]
\item \textbf{环物种。}美国加州中央谷地围着一圈暗色螈：山谷这一侧的邻居种群之间都能交配，基因像接力棒一样一圈圈传下去，可绕到山谷南端首尾相接时，两头碰面的种群却老死不相往来、无法杂交。这到底算一个物种还是两个？“生殖隔离”这把尺子在这里打了结。
\item \textbf{杂交区。}欧洲的小嘴乌鸦和冠鸦毛色一黑一灰，分明是两个物种，可在两者分布区交界的地带，它们稳定地杂交出一条窄窄的“混色带”，几万年来既不合并也不分家。物种界线在地图上是一段渐变区，而不是一条线。
\item \textbf{细菌根本不按这套规则玩。}它们不靠两性生殖，还能把基因直接“递”给毫不相干的邻居（水平基因转移）——抗生素抗性就这样在细菌之间传来传去。给细菌画的家谱更像一张反复打结的网，而不是一棵干净的树。这是第 79 章的主角，这里只提一句。
\end{itemize}

这些麻烦不说明“物种”概念没用——地图有边界，不代表地图没用。它们说明的是：\textbf{物种形成是一个还在进行中的连续过程，我们硬要在过程的任意瞬间贴标签，就总会遇到贴不稳的地方。}环物种恰恰是最珍贵的案例：它把“一个物种正在变成两个”的中间状态活生生摆在你眼前，相当于物种形成的慢动作回放。

\begin{tiaozhan}
为什么合子后隔离几乎总是雄性先不育（狮虎兽、骡子、果蝇杂种里都是雄性更严重）？这背后有一条叫“霍尔丹法则”的经验规律，遗传学家还在争论全部原因，但有一段机制已经很清楚，值得一看。设想祖先种群的一条染色体上有两个基因 A 和 B，它们要配合工作。种群分开后，一支把 A 升级成了 A'，另一支把 B 升级成了 B'——各自内部都运行良好。可当 A' 和 B' 在杂种体内相遇，这对从未一起测试过的“新零件”却合不来，发育就此卡壳。这就叫“多布然斯基–穆勒不相容”：\textbf{没有任何一个突变本身有害，伤害来自两个分开演化的基因在杂种里的第一次合作}。这个模型解释了一件反直觉的事——生殖隔离不需要谁“犯错”，它是分歧积累的副产品。至于为什么雄性更倒霉：雄性只有一条 X 染色体，X 上任何不合拍的基因都没有备份可遮掩（女性有两条 X，可以互相兜底）。细节超出主线，跳过不影响后面的阅读。
\end{tiaozhan}

\begin{shiyan}
\textbf{观察：两种“差不多”的动物，其实错开了生活。}

材料：笔记本、笔（可选：望远镜、手机拍照）。

步骤：
\begin{enumerate}[itemsep=1pt]
\item 在小区或公园里选定两种常见又容易混淆的鸟（例如麻雀和白头鹎）或两种相似的昆虫（例如菜粉蝶和东方菜粉蝶）。
\item 连续观察一周，每天早晚各一次，每次 15 分钟，记录：它们出现在什么位置（地面、灌木、树梢）？吃什么？一天中什么时段最活跃？叫声或飞行方式有何差别？
\item 周末把记录整理成一张对照表，回答：这两种动物有没有互相求偶或争抢同一处食物？它们的生活在哪些维度上“错开”了？
\end{enumerate}

观察点：把你发现的每一处“错开”对应到本章的隔离机制表——时间错开、空间错开、行为错开，各属于哪一类？

安全提示：只观察，不捕捉、不投喂、不靠近鸟巢；野外活动告知家长并结伴而行。
\end{shiyan}

\section{回到动物园那张海报}

现在可以回答开场的问题了。

狮虎兽的存在，恰恰证明狮子和老虎\textbf{已经是两个物种}，而不是推翻物种概念。理由有三：第一，在野外，狮子住在非洲草原和印度林地，老虎住在亚洲森林，地理上根本碰不上面，狮虎兽只是人工圈养硬凑出来的；第二，即使凑成了，合子后隔离立刻现身——雄性狮虎兽全部不育，说明两套基因组的配合已经出了故障；第三，狮子和老虎这两支已经分开演化了数百万年，各自的形态、捕猎方式、社会行为早已南辕北辙。狮虎兽不是“新物种”，它是两个物种之间残余亲缘的一次尴尬展示——就像两本页码已经错乱、但还能勉强互相引用几段的说明书。

而狗和狼的故事正好反过来：狗是大约一万多年前由狼驯化而来，分化时间太短，两套说明书还没来得及改乱，所以直到今天仍能互通有无——生物学家甚至把狗分类为灰狼的亚种，而不是独立物种。外表的巨大差异（吉娃娃对藏獒）只是少数“发育开关”基因被人工选择快速改写的戏剧化结果，第 79 章会告诉你，调控序列的微小改动足以掀起形态的大浪。

这套“看树不看梯子”的眼光，早已跳出动物园：

\begin{itemize}[itemsep=2pt]
\item \textbf{追踪传染病。}新冠疫情中，科学家把世界各地病毒样本的基因序列画成系统发育树，一眼看出哪个变异株从哪里传进来、在哪一带扩散——树成了流行病学家的地图。
\item \textbf{保护生物多样性。}保护经费有限，优先保谁？系统发育树帮保护生物学家找出“独自代表一整条古老枝系”的物种（比如鸭嘴兽），保住它们就保住了更大的一段进化历史。
\item \textbf{农业育种。}作物的野生近缘种是重要的基因仓库：知道谁和栽培种亲缘近、还能杂交，育种家才知道去哪片山坡找抗旱、抗病的基因。
\end{itemize}

不过，本章的树有一个心照不宣的假设：基因只从父母传给孩子，代代垂直向下。细菌递基因的案例已经暗示，这个假设在生命史上破过不止一个大洞——基因组会复制自己、会借用别人的片段、会把旧零件改装成新机器。这些“不垂直”的故事，是下一章的内容。

\section*{练一练}

\begin{sikaoti}
\item 画一棵包含人、黑猩猩、猕猴三个物种的系统发育树，并标出人和黑猩猩的最近共同祖先节点。然后把树在黑猩猩所在的节点处“旋转”一下重新画一遍，说明两张图表达的是同一棵演化关系树。
\item 某种蛙分布在一条大河两岸。西岸种群春天产卵，东岸种群夏天产卵。它们属于哪种生殖隔离？如果上游修了水坝使两岸水温趋于一致、产卵时间逐渐重合，你预测这两个群体的演化命运会怎样变化？说明理由。
\item 骡子不育，是因为它的 63 条染色体在减数分裂时无法正常配对。现有一种植物，两个近缘种杂交后杂种完全能育。仅从这一条信息，你能对这两个种的染色体和分化时间做出什么推测？
\item 环物种中，A 能和 B 交配，B 能和 C 交配，但 A 和 C 不能交配。按生物学物种概念，A、B、C 算几个物种？这个案例暴露了“生殖隔离”定义的什么困难？用一句话说清。
\item 一位同学说：“人是最高等的动物，因为人在进化树的最顶端。”请画出你理解的进化树形状，并用树的三个读图规则逐条反驳这句话。
\item 假设两个完全隔离的种群各自每代积累 70 个新突变，每代 25 年。估算它们分化 200 万年后，基因组上大约积累了多少个不同的突变位点（两边合计）？这个估算忽略了哪些会让真实差异变大或变小的因素？
\end{sikaoti}
